\pdfminorversion=4
\documentclass[aspectratio=169]{beamer}

\mode<presentation>
{
  \usetheme{default}
  \usecolortheme{default}
  \usefonttheme{default}
  \setbeamertemplate{navigation symbols}{}
  \setbeamertemplate{caption}[numbered]
  \setbeamertemplate{footline}[frame number]  % or "page number"
  \setbeamercolor{frametitle}{fg=white}
  \setbeamercolor{footline}{fg=black}
} 

\usepackage[english]{babel}
\usepackage[utf8x]{inputenc}
\usepackage{tikz}
\usepackage{courier}
\usepackage{array}
\usepackage{bold-extra}
\usepackage{minted}
\usepackage[thicklines]{cancel}
\usepackage{fancyvrb}

\xdefinecolor{dianablue}{rgb}{0.18,0.24,0.31}
\xdefinecolor{darkblue}{rgb}{0.1,0.1,0.7}
\xdefinecolor{darkgreen}{rgb}{0,0.5,0}
\xdefinecolor{darkgrey}{rgb}{0.35,0.35,0.35}
\xdefinecolor{darkorange}{rgb}{0.8,0.5,0}
\xdefinecolor{darkred}{rgb}{0.7,0,0}
\definecolor{darkgreen}{rgb}{0,0.6,0}
\definecolor{mauve}{rgb}{0.58,0,0.82}

\title[2020-11-20-cms-analysistools-workshop]{Survey and status of Pythonic HEP analysis tools}
\author{Jim Pivarski}
\institute{Princeton University -- IRIS-HEP}
\date{November 20, 2020}

\usetikzlibrary{shapes.callouts}

\begin{document}

\logo{\pgfputat{\pgfxy(0.11, 7.4)}{\pgfbox[right,base]{\tikz{\filldraw[fill=dianablue, draw=none] (0 cm, 0 cm) rectangle (50 cm, 1 cm);}\mbox{\hspace{-8 cm}\includegraphics[height=1 cm]{princeton-logo-long.png}\hspace{0.1 cm}\raisebox{0.1 cm}{\includegraphics[height=0.8 cm]{iris-hep-logo-long.png}}\hspace{0.1 cm}}}}}

\begin{frame}
  \titlepage
\end{frame}

\logo{\pgfputat{\pgfxy(0.11, 7.4)}{\pgfbox[right,base]{\tikz{\filldraw[fill=dianablue, draw=none] (0 cm, 0 cm) rectangle (50 cm, 1 cm);}\mbox{\hspace{-8 cm}\includegraphics[height=1 cm]{princeton-logo.png}\hspace{0.1 cm}\raisebox{0.1 cm}{\includegraphics[height=0.8 cm]{iris-hep-logo.png}}\hspace{0.1 cm}}}}}

% Uncomment these lines for an automatically generated outline.
%\begin{frame}{Outline}
%  \tableofcontents
%\end{frame}

% START START START START START START START START START START START START START

\begin{frame}{Needless to say, CMS is already using Python a lot!}
\vspace{0.5 cm}
\only<1>{\includegraphics[width=\linewidth]{01-github-cmssw-language.pdf}}\only<2>{\includegraphics[width=\linewidth]{05-github-anyroot-python-machinelearning-uproot.pdf}}
\end{frame}

\begin{frame}{But I'll be talking about the Scikit-HEP tools}
\large

\vspace{0.2 cm}
\hfill \includegraphics[height=3 cm]{scikit-hep-logo.pdf}

\vspace{-2 cm}
\begin{itemize}\setlength{\itemsep}{0.2 cm}
\item<1-> Experiment agnostic
\item<2-> Connects to industry-standard ecosystem (NumPy, etc.)
\item<3-> Provides capabilities that don't exist in the larger ecosystem

\vspace{0.2 cm}
\begin{itemize}\setlength{\itemsep}{0.2 cm}
\item<4-> \normalsize because it's domain-specific: e.g.\ Lorentz vectors
\item<5-> \normalsize because HEP is a pioneer: e.g.\ histograms-as-objects
\end{itemize}

\vspace{0 cm}
\item<6-> Technical: focuses on columnar data to mitigate Python's speed issue

\begin{uncoverenv}<6->
\begin{center}
\includegraphics[width=0.75\linewidth]{ease-vs-flexibility.pdf} \mbox{\hspace{1 cm}}
\end{center}
\end{uncoverenv}
\end{itemize}
\end{frame}

\begin{frame}{This is a survey talk}
\vspace{0.17 cm}
\begin{columns}
\column{1.15\linewidth}
\only<1>{\includegraphics[width=\linewidth]{survey-1.pdf}}\only<2>{\includegraphics[width=\linewidth]{survey-2.pdf}}\only<3>{\includegraphics[width=\linewidth]{survey-3.pdf}}\only<4>{\includegraphics[width=\linewidth]{survey-4.pdf}}
\end{columns}
\end{frame}

\begin{frame}{Things to know about: pyhf}
\large

\vspace{-0.2 cm}
\hfill \includegraphics[height=3 cm]{logo-pyhf.png}

\vspace{-2 cm}
CMS uses Combine and ATLAS uses HistFactory \\ for many-channel combined fits.

\vspace{0.25 cm}
pyhf is a Python implementation of HistFactory, \\ replacing RooFit with NumPy, TensorFlow, PyTorch, or JAX.

\vspace{0.25 cm}
\begin{uncoverenv}<2->
I'm mentioning this because the speedup is {\it dramatic} (even more so with GPUs).

\renewcommand{\arraystretch}{1.2}

\vspace{0.25 cm}
\begin{columns}
\column{0.4\linewidth}
SUSY sbottom analysis \hfill {\small (CPU)}

\vspace{0.1 cm}
\begin{tabular}{l r l}
\hline HistFactory & $>$ 50 & minutes \\
pyhf-NumPy & 121 & \textcolor{gray}{seconds} \\
pyhf-PyTorch & 31 & \textcolor{gray}{seconds} \\
pyhf-JAX & 9 & \textcolor{gray}{seconds} \\
\end{tabular}

\column{0.4\linewidth}

SUSY Wh 1$\ell$bb \hfill {\small (CPU)}

\vspace{0.1 cm}
\begin{tabular}{l r l}
\hline HistFactory & $>$ 10 & minutes \\
pyhf-NumPy & 39 & \textcolor{gray}{seconds} \\
pyhf-PyTorch & 7 & \textcolor{gray}{seconds} \\
pyhf-JAX & 4 & \textcolor{gray}{seconds} \\
\end{tabular}
\end{columns}

\vspace{0.2 cm}
\small
See \textcolor{blue}{\url{https://indico.cern.ch/event/938677}} (CMS B2G meeting)
\end{uncoverenv}

\begin{uncoverenv}<3->
\vspace{-3.9 cm}
\hfill \includegraphics[height=4 cm]{pyhf-gpu.png}
\end{uncoverenv}
\end{frame}

\begin{frame}{Things to know about: iminuit, zfit, mplhep, particle}
\large

\renewcommand{\arraystretch}{1.2}

\vspace{0.25 cm}
\begin{tabular}{p{0.2\linewidth} p{0 cm} p{0.7\linewidth}}
\uncover<1->{\includegraphics[width=\linewidth]{logo-iminuit.pdf}} & & \uncover<1->{Use MINUIT directly in Python with a pleasant (and newly revamped) interface.} \\

\uncover<2->{\includegraphics[width=\linewidth]{logo-zfit.png}} & & \uncover<2->{Aims to be the general purpose fitter with a Pythonic interface. Something to keep an eye on.} \\

\uncover<3->{\includegraphics[width=\linewidth]{logo-mplhep.png}} & & \uncover<3->{Make various HEP-standard plots in Matplotlib that conform to CMS or ATLAS style. \hspace{3.5 cm} \linebreak \linebreak Suspiciously familiar logo.} \\

\uncover<4->{\includegraphics[width=\linewidth]{logo-particle.pdf}} & & \uncover<4->{\vspace{0.25 cm}Pythonic access to PDG data. Don't hard-code masses!} \\

\end{tabular}
\end{frame}

\begin{frame}{Things to know about: Numba, Xarray, Zarr, Dask, Ray}
\large

\renewcommand{\arraystretch}{1.1}

\vspace{0.1 cm}
\begin{tabular}{p{0.1\linewidth} p{0 cm} p{0.8\linewidth}}
\uncover<1->{\includegraphics[width=1\linewidth]{logo-numba.png}\vspace{-0.4 cm}} & & \uncover<1->{JIT-compiler for Python: decorate a function with \mintinline{python}{@numba.jit} and it runs at the speed of C code. Some restrictions apply.} \\

\uncover<2->{\includegraphics[width=1\linewidth]{logo-xarray.png}\vspace{-0.4 cm}} & & \uncover<2->{Tracks metadata for the axes of an N-dimensional array, so you can slice it in natural units; built-in Zarr and Dask integration. \linebreak No jagged arrays---yet. (Awkward-Xarray integration?)} \\

\uncover<3->{\includegraphics[width=1\linewidth]{logo-zarr.png}\vspace{-0.4 cm}} & & \uncover<3->{Delivers slices of an array dataset to distributed jobs. \linebreak Also considering Awkward Array integration.} \\

\uncover<4->{\includegraphics[width=1\linewidth]{logo-dask.png}\vspace{-0.4 cm}} & & \uncover<4->{Distributed processing with (and without) an array interface. \linebreak Records lazy calculations and replays them on remote workers. \linebreak Also considering Awkward Array integration.} \\

\uncover<5->{\includegraphics[width=1\linewidth]{logo-ray.png}\vspace{-0.4 cm}} & & \uncover<5->{More robust distributed processing layer; can be used with Dask.} \\

\end{tabular}
\end{frame}

\begin{frame}{Short update on Uproot and Awkward Array}
\includegraphics[width=\linewidth]{uproot-awkward-timeline.pdf}

\begin{uncoverenv}<2->
\vspace{-4.28 cm}
\includegraphics[height=4.28 cm]{Raiders-of-the-Lost-Ark-Chamber.jpg}

\begin{center}
After Dec 1, ``\mintinline{python}{import uproot3 as uproot}'' if you need the old one.
\end{center}
\end{uncoverenv}
\end{frame}

\begin{frame}{Short update on pyjet and vector}
\large
\vspace{0.5 cm}
\begin{columns}
\column{0.3\linewidth}
\hfill \includegraphics[width=0.9\linewidth]{logo-pyjet.pdf}

\column{0.7\linewidth}
\textcolor{darkblue}{\Large Array-based interface to FastJet.}

\vspace{0.25 cm}
Currently passes pseudojets between Python and C++ as arrays, allows more efficient event-at-a-time access, rather than particle-at-a-time.

\vspace{0.25 cm}
Upgrade plans to pass jagged arrays of pseudojets and integrate with Awkward Array.

\textcolor{blue}{\small \url{https://github.com/scikit-hep/pyjet/issues/31}}

\end{columns}

\vspace{0.75 cm}
\begin{uncoverenv}<2->
\begin{columns}
\column{0.3\linewidth}
\hfill \includegraphics[width=0.9\linewidth]{logo-vector.pdf}

\column{0.7\linewidth}
\textcolor{darkblue}{\Large Array-based Lorentz vector calculations.}

\vspace{0.25 cm}
All the kinematic functions in NumPy, Awkward Array, Numba, and TensorFlow. Aiming for a minimal version by Dec 1 (excluding TensorFlow).

\end{columns}
\end{uncoverenv}
\end{frame}

\begin{frame}[fragile]{Short update on boost-histogram and hist}
\large
\vspace{0.5 cm}
\begin{columns}
\column{0.3\linewidth}
\hfill \includegraphics[width=0.9\linewidth]{logo-hist.pdf}

\column{0.7\linewidth}
\textcolor{darkblue}{\Large Histograms as accumulatable objects.}

\vspace{0.25 cm}
Histograms with a \mintinline{python}{fill} method are now a maintained part of Boost (C++ header-only library) and the Python interface is mature.

\end{columns}

\vspace{0.75 cm}
\begin{uncoverenv}<2->
\begin{columns}
\column{0.55\linewidth}
\includegraphics[width=\linewidth]{hist-docs.png}

\column{0.45\linewidth}
The \textcolor{darkblue}{hist} library wraps Boost Histogram in a ready-for-analysis way, bringing in all the plotting and statistical tools with

\vspace{0.25 cm}
\begin{minted}{python}
>>> import hist
\end{minted}

\vspace{0.25 cm}
Fully documented:

\textcolor{blue}{\small \url{https://hist.readthedocs.io/}}
\end{columns}
\end{uncoverenv}
\end{frame}

\begin{frame}{}
\Huge
\vspace{1 cm}
\begin{center}
\textcolor{darkblue}{Questions?}
\end{center}
\end{frame}

\begin{frame}{Next up:}
\Huge
\vspace{0.75 cm}
\begin{center}
\includegraphics[width=0.35\linewidth]{logo-coffea.pdf}
\end{center}
\end{frame}

\end{document}
